\scanpage{74}
\deleted{ries} $\psi^\circ:\overset{\substack{\inserted{S'=R'^\circ}\\\Vert}}{T'(\mathrm{Ens})}\longrightarrow\overset{\substack{\inserted{S=R^\circ}\\\Vert}}{T(\mathrm{Ens})}$ est simplement le foncteur associé à $T'\longrightarrow T$ déduit de $\varphi$. Ce foncteur doit donc commuter aux $\varprojlim$ quelc. [\unsure{de ce|ce|donc ce} qui exprime simplement que $\varphi$ commute aux $\varinjlim$, \inserted{\deleted{\unsure{et on|alors}} or c'est un adjoint à g.}] Donc les sous-théories de $T$ correspondent aux sous-catégories pleines \inserted{$S'$} de $S$, \deleted{\unsure{admettant un|et permettant de|admettant des}} \inserted{[stables par} $\varprojlim$, ce sera automatique par les autres conditions] admettant des $\varinjlim$ quelconques (mais pas néc. stables par ces dern) et telles que le foncteur d'inclusion $S'\overset{i\,(=\psi^\circ)}{\hookrightarrow}\unsure{S|S^\circ}$ ait un adjoint à \uline{gauche} \deleted{$\unsure{\varphi}$} $j\ (= \varphi^\circ)$
\[
\mathrm{Hom}_{S}(x,i(y'))\simeq\mathrm{Hom}(jx,y')
\]

\deleted{Alors on a des foncteurs, \illegible{1} adjts de \unsure{oubli que|oubliques}} \deleted{[Mais en général on (Structure} ($S'$ (foncteur $S'$-structure ``libre'' définie par une $S$-structure) Mais en général, il n'y a pas moyen de définir pour tte $C\in\mathrm{Ob}\,\Delta$ un foncteur $T(C)\longrightarrow T'(C)$ qui \deleted{\unsure{en|les}} redonne $j$ \deleted{\unsure{comme $C$}} pour $C=(\mathrm{Ens})$ [\uline{NB} Cependant, si $C$ est un \uline{topos}, et si $T,T'$ et $\varphi$ se définissent au stade des $\varprojlim$ finies, c'est Ok \ldots]

\uline{NB} La 2-catégorie des théories algébriques pour $\Delta$ est 2-équivalente à la 2-catégorie des cat. $S$ qui sont stables par \inserted{petites} $\varinjlim$ \deleted{quelc.} et qui ont une petite sous-catégorie str. génératrice, avec les foncteurs qui commutent aux $\varinjlim$ (mais en \unsure{renversant|inversant} les flèches pour les morphismes ?). \deleted{Ce sont aussi les catégories telles que les \inserted{contrav.} foncteurs qui transforment $\varinjlim$ quelc en limites \unsure{proj.}}
% archive sheet number: 73
